% Separate post-CP37 exploratory candidate; not wired into the frozen paper.
\subsection{An actual scalar element excludes the quadratic Cartan exception}

\begin{lemma}[Two incompatible coordinates for an order-$p^2$ automorphism]
\label{lem:B-RD102-scalar}
\status{Candidate; complete proof awaiting independent audits.}
Let $p\ge5$, $q=p^2$, and let $F/S$ be a totally
ramified finite Galois extension of complete
mixed-characteristic $(0,p)$ discrete valuation fields
with algebraically closed residue field.
Suppose $g\in\operatorname{Gal}(F/S)$ has order $p^2$
and $x_1,x_2\in\mathfrak m_F$ satisfy, for the normalized
valuation $v_F$,
\[
\begin{gathered}
 v_F(x_i)=A_i,\qquad v_p(A_i)=1,\qquad A_1>pA_2,\\
 v_F(gx_i-x_i)=qA_i,\qquad
 v_F(g^px_i-x_i)=q^2A_i,\qquad
 q^2A_1<v_F(p).
\end{gathered}                                           \tag{B-RD102.1}
\]
Then these assumptions are contradictory.
\end{lemma}
\begin{proof}
Choose a uniformizer of $F$ and a coefficient ring
inside $\mathcal O_S$, so the residue representatives
are fixed by $g$. Write
$a=b(g)=v_F(g\pi-\pi)-1$ and
$b=b(g^p)$. Both are positive.
The first ramification coefficient is additive on
each lower graded group, so taking the $p$th power
kills it. Since $g^p\ne1$, this gives $b>a$.
The full uniformizer expansion also gives the
unconditional inequality
$v_F(gx-x)\ge v_F(x)+a$.
Applied to $x_2$, it implies
\[
                a\le(q-1)A_2<\frac{(q-1)A_1}{p}.
                                                        \tag{B-RD102.2}
\]
Thus the critical equality for the first coordinate
is strictly avoided, rather than assigned a limiting case.

Use the fixed complete expansion of $x_1$.
Below $v_F(p)$, the term of degree $A_1$, which has
exact $p$-adic order one, first changes in degree
$A_1+pa$. Every later exponent divisible by $p$
changes strictly later. If $J$ is the first nonzero
ordinary exponent, its change is in degree $J+a$.
This follows directly by reducing the binomial expansions
below the stated cutoff, with fixed coefficients, exactly
as in the complete B60 and B81 coefficient proofs.
By (B-RD102.2), $A_1+pa<qA_1$, whereas the actual
change has degree $qA_1$. Cancellation is therefore
necessary and forces
\[
                     J=A_1+(p-1)a,\qquad p\nmid J ,
                                                        \tag{B-RD102.3}
\]
with a nonzero coefficient. If this asserted degree
were divisible by $p$, or its coefficient vanished,
the unique earliest change could not cancel, already
contradicting the actual displacement.

Now apply the same expansion to $g^p$, of break $b>a$.
The fixed ordinary term is uniquely first because
\[
             A_1+pb-(J+b)=(p-1)(b-a)>0 .
                                                        \tag{B-RD102.4}
\]
All later terms change later. If the first changed
degree were at least $v_F(p)$, every change would be
above the actual terminal degree, which is impossible.
Thus the ordinary term is below the cutoff, and its
nonzero coefficient cannot cancel. The exact terminal
displacement gives
$J+b=q^2A_1$, that is
$b=(q^2-1)A_1-(p-1)a$.
The unconditional inequality for $g^p$ on $x_2$
gives $b\le(q^2-1)A_2$. Combining it with (B-RD102.2)
would imply
\[
 (q^2-1)(A_1-A_2)\le(p-1)(q-1)A_2 .
                                                        \tag{B-RD102.5}
\]
But $A_1>pA_2$ makes the left side greater than
$(p-1)(q^2-1)A_2$, which is strictly larger than
the right side. This contradiction proves the lemma.
Only the two actual elements $g,g^p$ and their
fixed coordinate expansions were used.
\end{proof}

\begin{proposition}[The quadratic tame exception has more than four matrix dimensions]
\label{prop:B-RD102-tame}
\status{Candidate; complete proof awaiting independent audits.}
Let $K$ be complete discretely valued of mixed
characteristic $(0,p)$, $p\ge5$, with algebraically
closed residue field, $\mu_p\subset K$ and
$e_K=p(p-1)$. A pair of whole potentially supersingular
height-two groups with isomorphic actual $p^3$-torsion,
tame reducible residual inertia, and different finite
special characters has wild division image of order
at least $p^5$.
For $p\equiv1\pmod4$, the already proved stronger
bound $p^6$ remains valid.
\end{proposition}
\begin{proof}
Let $N=p^2-1$ and use the actual common tame good
field $L/K$ of degree $N$.
B77 excludes square common finite determinant.
The complete actual second-point comparison
Corollary~\ref{cor:B-RD84-first-wild} forces
\[
 B_i=p^2k_i,\quad1\le k_i\le p-2,\qquad
 H_1=\mathbf F_pe_1,\quad H_2=\mathbf F_pe_2 .
                                                        \tag{B-RD102.6}
\]
In particular both first-level modules are $L$-rational.
The two actual quotient lattice computations in B83
force nonzero opposite off-diagonal projections of
the same first layer $W_1$. Lemma~\ref{lem:B-RD85-diagonal}
forces both diagonal projections to be nonzero.
Actual matrix powers and commutators give
$W_1\subset W_2$ and $[W_1,W_1]\subset W_2$.

Let the tame residual ratio be $\rho$.
It is nonsquare, so $\rho\ne1$. If $\rho\ne-1$,
the three weights $1,\rho,\rho^{-1}$ are distinct.
Their integral polynomial projectors put both
$E_{12},E_{21}$ and a nonzero diagonal direction
in $W_1$. Thus $\dim W_1\ge3$, and the image has
order at least $p^6$.
It remains $\rho=-1$, necessarily with
$p\equiv3\pmod4$. The diagonal and off-diagonal
weight spaces still split. Hence $\dim W_1\ge2$.
If the image had order at most $p^4$, the actual
order formula and $W_1\subset W_2$ would force
$\dim W_1=\dim W_2=2$.

Write $W_1=\langle D,A\rangle$, with $D$ diagonal,
both its entries nonzero, and $A$ off-diagonal,
both its entries nonzero. Since $W_2=W_1$,
the commutator $[D,A]$ must be a scalar multiple
of $A$. In odd characteristic this forces the two
entries of $D$ to be equal: otherwise the opposite
signs in the two off-diagonal positions give a second
independent off-diagonal direction.
Thus $I\in W_1$. There is an actual element
$g=1+pI+p^2C$ in the third-level division group,
and
\[
                     g^p=1+p^2I\pmod{p^3}.
                                                        \tag{B-RD102.7}
\]
It has exact order $p^2$.

The high-line exponent in model $i$ is $c_i-B_i$;
the other residual exponent is $pc_i+B_i$.
The ratio $-1$ implies
$(p-1)c_i+2B_i\equiv N/2\pmod N$.
Reduction modulo $p-1$ gives $2k_i\equiv0\pmod{p-1}$.
The range in (B-RD102.6) therefore forces
$k_1=k_2=(p-1)/2$. Their high and low values are
\[
 \alpha=\frac{p(p-1)(p+2)}2,\qquad
 \beta=\frac{p(p-1)}2,\qquad \alpha=(p+2)\beta .
                                                        \tag{B-RD102.8}
\]

Let $F=L(M_i[p^3])$, which under the contrary
assumption has degree $p^4=q^2$ over $L$.
Choose a corresponding primitive third-level point
$P_i$ whose actual first-level image $p^2P_i$
is $e_1$. It is high in model1 and low in model2.
Since $B=p^2(p-1)/2>e$, the full inverse Newton
polygon over a high target is a single $q$-segment;
the next inverse also is, as its target is still smaller.
The same is true for both low inverse steps.
Thus the normalized $F$-values of the two point
coordinates are $A_1=\alpha$ and $A_2=\beta$.
Both have exact $p$-adic order one.

The actual displacement $(g-1)P_i$ has order $p^2$
and its $p$-multiple is $p^2P_i=e_1$, since the
$p^2C$ term is killed by this multiplication.
Its valuation in $F$ is consequently $qA_i$.
Equation (B-RD102.7) gives terminal displacement
$p^2P_i$, of value $q^2A_i$.
Formal subtraction differs from coordinate subtraction
by an integral unit. All these degrees lie below
$v_F(p)=eNq^2$, since $\alpha<eN$.
Finally $A_1=(p+2)A_2>pA_2$.
The preceding actual-automorphism lemma now gives
a contradiction. Thus this quadratic Cartan exception
also cannot have order at most $p^4$.
\end{proof}

\begin{remark}[Scope of the actual exclusion]
\label{rem:B-RD102-scope}
\status{Candidate; no larger image is excluded or constructed.}
This is a finite-level actual-field argument, not an
inference of characteristic-zero CM from a Cartan
congruence image. It closes only the order-at-most-$p^4$
tame reducible exception when $p\equiv3\pmod4$.
It does not settle the full first-wild problem at
$p=7$ or any larger remaining image.
The fields, the order-$p^2$ element and its
$p$th power, and the two coordinates are all inside
the same actual division field.
No equality of breaks is asserted across unidentified
quotients. The uniform proof uses the complete B60 and
B81 fixed-coefficient arguments, the B69 actual inverse
polygon, and the B83--B85 actual quotient and point-field
proofs. These dependency proofs were reread; no fresh
external-source opening is claimed here.

\smallskip
\noindent\emph{Exact finite supplement.}
The retained program
\path{code/reduction_descent/verify_tame_quadratic_Cartan_exclusion.py}
checks six symbolic identities, 1,108 character rows
for the 24 primes from 5 through 101, the 12 quadratic
rows and their 48 inverse-polygon comparisons, 17,667
literal matrix $p$th powers and displacement checks
modulo $p^3$ for $p=5,7,11$, 11,552 commutator-plane
checks, and 720 fixed-expansion cases using 285 direct
binomial computations. The author run passed in
$0.5161670989909908$ seconds. No execution failed.
These finite checks neither realize a Galois field nor
replace the actual-field proof. The exact executed
program, output and prior candidate remain retained.
Independent full audits remain required before promotion.
\end{remark}
